\chapter{第二十五讲\quad 危害国家安全罪}
\label{lecture:25}
\noindent{\small\hyperref[ref:lecture:25]{查阅本讲扩展规则}}\par

教材本讲集中讨论第110、111、109条。把「组织任务」「境外对象与资料」「身份及叛逃行为」分作三条证明链，每条单独核对主客观事实。

\subsection*{一、间谍罪}

\paragraph{法条与定义}
\(J\) 表示参加间谍组织，\(T\) 表示接受间谍组织及其代理人的任务，\(B\) 表示为敌人指示轰击目标，\(H\) 表示危害国家安全。第110条的条文结构为
\[
(J\lor T\lor B)\land H\ \Longrightarrow\ B_{110},
\]
并须核对行为人的认识和故意。

\paragraph{推导与边界}
甲明知对方为间谍组织代理人，接受其有危害国家安全的任务，证据足以认定 \(T,H\) 和相应故意。即使未另行证明 \(J\)，析取式已有真项，可依第110条进行评价。普通境外机构发出的合法合作邀请，未证明组织或任务属性，不能推出 \(T\)。是否另有获取国家秘密等行为，应另列事实、另查罪数。

\subsection*{二、为境外窃取、刺探、收买、非法提供国家秘密、情报罪}

\paragraph{法条与定义}
令 \(W\) 为行为所服务的是境外机构、组织、人员，\(S\) 为法定国家秘密或情报，\(A\) 为窃取、刺探、收买、非法提供之一，\(K\) 为相应故意。第111条可写为 \(W\land S\land A\land K\Rightarrow B_{111}\)。罪名包含数个可选行为，最终称谓依认定事实确定。必须先查资料性质，再查对象以及行为人对有关事实的认识。

\paragraph{推导与边界}
甲明知某机构在境外，为其非法交付经证实属于国家秘密的文件。收件方、文书性质、交付行为及沟通记录分别给出 \(W,S,A,K\)，依第111条推出相应非法提供分支。若甲确实不知道对象在境外，不能保留同一故意谓词，须另核第398条故意泄露国家秘密及情节等要求。若只有以窃取、刺探、收买方法非法获取国家秘密，并无为境外的联系，另查第282条。公开可得的一般信息不能仅凭接收方位于境外即认作法定情报。

\subsection*{三、叛逃罪}

\paragraph{法条与定义}
令 \(I\) 为国家机关工作人员，\(D\) 为履行公务期间，\(L\) 为擅离岗位，\(E\) 为叛逃境外或在境外叛逃。第109条第一款核对 \(I\land D\land L\land E\)。第二款另明定掌握国家秘密的国家工作人员叛逃境外或在境外叛逃，依照第一款的规定从重处罚。第二款的主体和规定不能简单改写成第一款每项条件的重复。

\paragraph{推导与边界}
国家机关工作人员甲在境外履行公务时，故意擅离岗位并实施经查明的叛逃行为。由身份、公务安排、擅离事实及叛逃证据依次得 \(I,D,L,E\)，根据第一款评价。若甲只是依法离职后获准出境，已证实无叛逃，则 \(E\) 为假，第一款无从推出。若甲叛逃完成后又接受间谍组织任务，应把后行为单列，按教材分别核对两罪和并罚。

\paragraph{闭卷记忆}
\begin{enumerate}
\item 写三组门槛：「任务、资料、离岗」。任务组补「间谍组织、代理人或为敌指示目标」；资料组补「为境外、秘密或情报、四个动作」；离岗组补「身份、公务期间、擅离、叛逃」。
\item 每组再问一次「明知什么」。末尾单独比较同一行为与先后两个行为，避免把境外联系当成间谍任务的证明。
\end{enumerate}
口诀「\textbf{任务查来源，资料查对象，叛逃查职务与行为}」。

\paragraph{闭卷自测与答案}
\begin{enumerate}
\item 甲接受明知为间谍组织代理人的任务，未正式入会，是否因此排除第110条？答：不能。接受任务是独立析取项，仍须证明危害国家安全及故意等要件。
\item 甲向境外机构发去资料，证明资料属于国家秘密之前，能否只凭收件地址推出第111条？答：不能。\(W\) 即使已证实，\(S\) 仍须证明，也要核对法定行为与故意。
\end{enumerate}
